\begin{abstract}
\noindent
AI agents now pause for approval before they send money, delete files or email outsiders, and the approval,
whether given by a person or by an automated approver, is usually treated as a second factor. We argue that it
is a second factor only when compromising the agent cannot produce it, and only when the approval says exactly
what it approves, and we test both conditions. We define twelve attacks on the approval path, among them decoy
recipients, amount cancellation, card padding, label spoofing, enrollment squatting and threshold splitting. We
execute them against the Agent Payments Protocol (AP2), Microsoft's action-bound approval design and Model Context
Protocol (MCP) elicitation, review OpenClaw statically, and map them onto agent frameworks through an audit of ten
agent runtimes. By our reading of their APIs, no runtime produces a third-party verifiable approval record. None
of the seven runtimes where we could measure flooding limits approval floods within a run, and in four of the
eight whose binding we could measure, a resumer, a hook or a stored history can change what executes after
approval. As a reference defense we build \sys{} (Agent Step-Up), a gateway that canonicalizes each tool call
into an effect descriptor, renders the approval card from that descriptor, has a passkey on the user's phone sign
its digest, and executes only the stored descriptor. Under three stated hypotheses, Tamarin proves its binding
and single-use properties and key-level non-repudiation, and finds the expected attacks on four baselines, a
mutant and a straw man (\foLemmas{} lemmas, all with the expected verdict). Across \scTrials{} scripted trials
and \adPairs{} AgentDojo attack pairs, \sys{} stops substitution, replay and grant stretching, which unbound
approval allows every time. Against an approver that compares the card with the task, the trusted card cuts
forged-narration attack success from \adFExNarr{} to \adFExAsu{}. With \llmModels{} current language models
acting as automated approvers, the trusted card raises the attack catch rate from \llmCatchNarr{} to
\llmCatchCard{}, at the cost of refusing \llmFDPctCard\% of harmless near misses, most of them because of a
placeholder artifact in our stimuli. We also report what does not work: the second device adds nothing measurable
over a trusted on-host confirmation in these experiments, flooding still succeeds against \sys{}
\scFloodPthirty{} to \scFloodDefault{} of the time with a modeled human approver, and approval cannot stop effects
that policy already allows. Code, data and proofs are available at \codeurl{}.
\end{abstract}

\noindent\textit{Keywords}\enspace AI agents $\cdot$ step-up authorization $\cdot$ multi-factor authentication
$\cdot$ WebAuthn $\cdot$ effect binding $\cdot$ trusted rendering $\cdot$ formal verification

\bigskip
